% =============================================================================
\chapter{Limitações e trabalho futuro}
\label{cap:limitacoes}
% =============================================================================

Este capítulo reúne o que o sistema não faz hoje, cada item com o motivo e,
quando há, o caminho para resolvê-lo. Nada aqui impede o uso em aula; são os
pontos de partida para quem continuar o trabalho.

\section{Hardware}

\paragraph{Tamanho fixo de 1024 amostras.} Os aceleradores foram sintetizados
para blocos de 1024. Sinais maiores são tratados no cliente, por blocos
(seção~\ref{sec:blocos}), com o custo de uma requisição por bloco. Aumentar o
bloco exige memórias maiores, um núcleo de FFT gerado de novo e, para a
convolução, aceitar um tempo que cresce com o quadrado do tamanho: 1024 por
1024 já leva 170\,ms.

\paragraph{A licença de avaliação da FFT.} O bitstream só funciona com o
\emph{tether} vivo, e para uma hora depois de ele cair
(seção~\ref{sec:tether}). O \arq{morphe-up.sh} torna isso invisível na rotina,
mas a estação precisa ficar ligada durante a aula e o \arq{morphe-up.sh} tem de
rodar depois de cada boot dela. As saídas são uma licença do núcleo da Intel,
ou uma FFT própria em Verilog, que tiraria a restrição e daria controle sobre a
aritmética interna (hoje o teto de 92\,dB é do núcleo).

\paragraph{A aritmética do \texttt{conv1d}.} A convolução e o FIR truncam os
produtos e deixam o acumulador dar a volta. Trocar por arredondamento e
saturação reduziria o erro máximo de 521 para 25 LSB e eliminaria o estouro
silencioso. A troca foi avaliada e deliberadamente não feita, para manter o
comportamento descrito no TCC original (seção~\ref{sec:acel-conv}); é uma
decisão a rever com o supervisor, não uma pendência técnica.

\paragraph{Restos do projeto original.} O IP de FIR da Intel (\texttt{fir\_ii})
e o seu adaptador continuam no projeto sem instância, e o registrador
\texttt{fir\_error} não é acionado por nada. Podem sair numa próxima
compilação, junto com a conferência de que nada mudou.

\section{Aquisição}

\paragraph{O que falta validar.} Na data deste manual o ADC foi validado na
placa com tensões contínuas (0\,V e 3,3\,V, contra o valor nominal) e com a
captura contínua a 200\,kHz sem perda; \textbf{falta a validação com uma
senoide do gerador de funções} (SINAD e ENOB medidos, e a continuidade das
emendas entre blocos com um sinal que muda) e a comparação da leitura em 3,3\,V
com um multímetro. O roteiro está em \arq{docs/COMPILAR-ADC.md}, passos 7c e
7d, e no \arq{testa_adc.py}. Com a mesma senoide falta conferir na placa o modo
\emph{ao vivo} da janela de aquisição e o analisador espectral, que em
30/09/2026 foi testado na placa só com a entrada sem sinal.

\paragraph{Um transitório no início da captura.} Com a entrada do canal~0 da
placa~1 solta, toda captura começa perto de 2,4\,V e decai, em 40 a 50
amostras, até um nível estável que depende da taxa (1,96\,V a 10\,kHz, 1,84\,V a
50\,kHz). Parece efeito da entrada sem nada ligado, de alta impedância, que o
próprio conversor carrega. Se o transitório continuar com o gerador ligado, é
defeito da captura, afeta também a janela de aquisição, e precisa ser
investigado.

\paragraph{Não há filtro \emph{anti-aliasing}.} O sinal chega ao conversor sem
filtro analógico. Um tom acima de \(f_s/2\) aparece dobrado no espectro. Para
medir, isso é uma limitação; para ensinar amostragem, é uma oportunidade.

\paragraph{Faixa de entrada.} A entrada lê de 0 a 4,095\,V e pode ser danificada
abaixo de \(-0{,}3\)\,V. Um sinal alternado precisa de deslocamento DC antes do
conector. Um pequeno circuito de condicionamento (deslocamento, proteção e
filtro) tornaria a entrada segura para uso direto pelos alunos.

\paragraph{Um canal por captura.} A captura lê um canal, ou um par
diferencial, por vez. Como o LTC2308 recebe a palavra de configuração a cada
conversão, alternar dois canais é uma mudança pequena no \arq{adc_captura.v}
(e exige recompilar); abriria correlação cruzada e estimativa de atraso.

\paragraph{A placa fica ocupada durante a captura.} A captura contínua mantém
a conexão aberta, e o servidor não atende mais ninguém nesse tempo
(capítulo~\ref{cap:placas}). Limitar a duração, ou atender outras operações
durante a captura, são decisões ainda em aberto.

\section{Infraestrutura}

Descritas no capítulo~\ref{cap:placas}. Três melhorias resolveram a maior
parte do que faltava: a \textbf{porta de estado} (28/09/2026) diz a qualquer
momento o que cada placa está fazendo, e o cliente escolhe a placa livre e menos
usada; o \textbf{\arq{placas.conf}} versionado (28/09/2026) dá a todos os
computadores a mesma lista de placas; e a \textbf{busca por \emph{broadcast}}
(30/09/2026) acha as placas pelo nome, no IP que tiverem, o que tornou
desnecessária a reserva de DHCP. O que continua faltando: \textbf{reserva} de
placa (a escolha divide a turma, mas não impede duas pessoas na mesma placa) e
\textbf{troca automática de placa} durante o uso. Falta também testar dois
computadores abrindo o aplicativo ao mesmo tempo: devem cair em placas
diferentes.

A mesma porta de estado protege quem está usando uma placa: o
\arq{morphe-up.sh} a consulta antes de reprogramar e pergunta se deve
continuar quando a placa está em uso. Mas a pergunta só funciona para quem roda
o script; nada impede, por exemplo, alguém de gravar um projeto próprio pelo
Quartus numa placa que está servindo a turma.

\paragraph{Atualização dos computadores.} Em 28/09/2026 o Quartus e o
aplicativo estavam em todos os computadores, cada um instalado pelo seu script
(\arq{instala-quartus.sh} e \arq{instala-morphe.sh}). Mas o aplicativo não se
atualiza sozinho: desde 01/10/2026 basta um comando por computador,
\arq{/opt/morphe/atualizar.sh} (seção~\ref{sec:atualizar}), mas ainda é preciso
lembrar de rodá-lo em cada um, e com vários computadores isso acaba esquecido.
A proposta é um temporizador do sistema (\emph{systemd timer}),
instalado pelo próprio \arq{instala-morphe.sh}, que atualize o
\arq{/opt/morphe} como a conta dona, ao ligar o computador e a cada hora, só
por avanço simples (\opt{ff-only}): um arquivo modificado à mão continua
protegido, porque nesse caso a atualização é recusada.

\paragraph{A estação.} Três coisas dependem dela: o \emph{tether} da licença
(ela precisa ficar ligada durante a aula, e o \arq{morphe-up.sh} tem de rodar
depois de cada boot), os cabos JTAG das placas, e a compilação do bitstream. O
servidor SSH que ela ganhou em 23/09/2026 só servia para os outros computadores
copiarem o Quartus, e foi desligado em 30/09/2026.

\section{Ideias que a plataforma já permite}

Com o que existe hoje, sem mudar o hardware:

\begin{itemize}
  \item aula de amostragem e \emph{aliasing}, capturando um tom do gerador
    acima e abaixo de \(f_s/2\);
  \item com o analisador espectral (seção~\ref{sec:analisador}): o espelho em
    \(f_s - f\), o efeito da janela na amplitude de um tom, a resolução
    \(f_s/N\) separando dois tons próximos;
  \item medir o próprio conversor (SINAD, ENOB) e comparar com os 74\,dB
    teóricos de 12 bits;
  \item multitaxa: capturar rápido e dizimar com o FIR da placa;
  \item os exercícios do Oppenheim que usam sequências dadas, agora que o
    aplicativo lê sinais de arquivo de qualquer tamanho;
  \item estudar o efeito da quantização dos coeficientes de um IIR de ordem alta
    em forma direta contra seções de segunda ordem, com o projetista e o
    comparador.
\end{itemize}
